\section{Appendix: the max-smoothing lemma}

\DMGtodo[inline]{This was just copy pasted from the old preprint on Dehn-Thurston. Do we need to reconcile the language in this appendix with the rest of the paper? Or should we keep them separate? }

\subsection{Preliminary language on taut curves}

The following has been extracted from the unpublished preprint~\cite{Thu00:GeometricIntersection}.
The context of the proof is more general than what we use in our main theorem, since it particular it allows surfaces with boundary and arc components between them. For the purposes of the proof, it is useful to distinguish the taut representatives in a homotopy class of curve, also. We keep the same notation as in that preprint for now, so we warn the reader that this section uses slightly different terminology from that in the rest of the paper.

\begin{definition}
  \label{def:curve}
  A \emph{curve} on $\Sigma$ is a properly immersed 1-manifold, possibly
  with boundary.  The set of all curves on $\Sigma$ is denoted $\C(\Sigma)$.
  $\C_0(\Sigma)$ is the set of curves with no boundary, and for $S$ a
  subset of the boundary components of $\Sigma$, $\C_S(\Sigma)$ is the set of
  curves with boundary only on~$S$.

  The circle components of~$C \in \C(\Sigma)$ are called \emph{loops} and the
  other components are called \emph{arcs}.
\end{definition}

Note that a curve is what other authors have called a multi-curve or
sometimes a curve system.

\begin{definition}
  \label{def:taut}
  A curve $C \in \C_S(\Sigma)$ is \emph{taut} if it is immersed with the
  minimum number of self-intersections in its homotopy class, it has
  no trivial components, and it has no components which are homotopic
  into any of the boundary curves in~$S$.  If $C$ is not taut, it is
  said to have \emph{excess intersections}.  Let $\TC_S(\Sigma)
  \subset \C_S(\Sigma)$ be the set of taut curves.
%\nota{excess intersections is bad name}
%\nota{Use $\TC$ notation}
\end{definition}

\DMGtodo[inline]{This is the main different piece of language in this section, the notion of tautness. We don't use it previously, but we could define it here and use the same proof.}

We exclude these various trivial cases so that for any $C \in
\TC_S(\Sigma)$, there
is a $D \in \TC_{\bar S}(\Sigma)$ which intersects $C$ essentially, where
$\bar S$ is the set of connected components of $\partial \Sigma \setminus S$.

There is a nice geometric characterization of tautness.

\begin{theorem}[Freedman-Hass-Scott~\cite{FHS82:ClosedGeodesics},
  Neumann-Coto~\cite{NC01:ShortestGeodesics}]
  \label{thm:taut-minimizing}
  Let $C$ be a curve with no components which are proper powers of
  other curves and no two parallel components.
  Then $C$ is taut if and only if it is length-minimizing in its
  homotopy class with respect to some Riemannian metric on~$\Sigma$.
\end{theorem}
%\nota{Should I also cite Marc
 % Shepard~\cite{Shepard91:TopologyShortest}?  Do I need to say
 % anything about the boundary?}
  
  

(The reason for the restriction on $C$ in
Theorem~\ref{thm:taut-minimizing} is that in the excluded cases,
length-minimizers are covers of a single curve and do not intersect
themselves transversally.)
\DMGtodo[inline]{This is in contrast to our definition of essential crossing in From curves to currents, where we do consider this ``self-intersections'' of a non-primitive curve as essential.}

\begin{corollary}
  \label{cor:taut-extension}
  If $C\in \C(\Sigma)$ is tautly embedded and $D\in \C(\Sigma)$ is arbitrary, then
  there is a taut embedding of $C \cup D$ which agrees with the given
  embedding of $C$.
\end{corollary}%\nota{Vague about exact boundary conditions and
  %tautness.  Not quite clear how to most elegantly deal with multiple
  %components.}

We will generally be interested in parameterizing simple closed
(multi-)curves.

\begin{definition}
  \label{def:simple-curve}
  A \emph{simple curve} on~$\Sigma$ is an taut, embedded closed
  curve~$C$, not necessarily connected.
  The space of simple curves up to ambient isotopy (or equivalently up
  to homotopy) is denoted $\SC(\Sigma)$.  Similarly for $\SC_0(\Sigma)$ and
  $\SC_S(\Sigma)$.
\end{definition}
\DMGtodo[inline]{In the paper we say simple multi-curves (no weights, here)}
Our ``simple curves'' are also called ``curve systems'' or ``integer
measured laminations''.  The name ``integer measured lamination''
comes from the general theory of measured laminations; all results in
this paper extend to the more general case.

To parametrize simple curves, we will use intersection numbers.
\begin{definition}\label{def:intersection}
  Given $C\in \C_S(\Sigma)$ and $D\in \C_T(\Sigma)$, where $S$ and~$T$ are
  disjoint, the \emph{geometric intersection number}~$i(C,D)$ is the number of
  intersections between $C$ and $D$ in a taut embedding of $C \cup D$.
  We will also write this as $i_C(D)$ or $i_D(C)$.
\end{definition}
%\nota{I use more than this for the twisting case.  Have to deal
 % properly with issues of tautness.}
  
 \subsection{The max-smoothing lemma}
 
 The following lemma is a central tool in this paper.

\begin{lemma}[Max-smoothing lemma]\label{lem:smoothing}
  Let $S$ and $T$ be disjoint subsets of $\partial \Sigma$, \DMGtodo[inline]{In the paper we don't consider $\Sigma$ with boundary or punctures.} and let $X \in \SC_S(\Sigma)$
  and $A \in \C_T(\Sigma)$, with $A$ tautly embedded.  Let~$x$ be a crossing of \DMGtodo[inline]{Recoincile with essential crossing. For self-intersections of non-primitive curves there is a discrepancy.}
  $A$, and let $A_+$ and $A_-$ be the two smoothings of~$A$ at~$x$.  Then
  \[
  i_X(A) = \max(i_X(A_+), i_X(A_-)).
  \]
  \label{lem:maxsmoothing}
\end{lemma}

For instance, suppose $A$, $B$, $C$, $D$, $E$, and~$F$ are the 4 sides
and two diagonal of a quadrilateral in~$\Sigma$.  (The corners are four
boundary components of~$\Sigma$, possibly with repetitions.)  We can use
this theorem to compute $i_X(F)$ in terms of the number intersections
of $F$ with the other 5 sides: we consider the non-simple curve $E \cup
F$.  This curve has one self intersection; if we smooth out the
intersection, we get two opposite sides of the square, so we have
\begin{align*}
i_X(E \cup F) &= \max(i_X(A \cup C),i_X(B \cup D))\\
i_X E + i_X F &= \max(i_X A + i_X C, i_X B + i_X D).
\end{align*}
If we use the number of intersections with the edges of a
triangulation to parametrize the different curves~$X$, this equation
lets us change coordinates from one triangulation to another:
\[
\mathgraph{moves-0} \mapsto \mathgraph{moves-1}
\]
This special case of Lemma~\ref{lem:smoothing} was known before
and is in any case easy to derive by a case analysis.

Before we prove Lemma~\ref{lem:smoothing}, let us study excess
intersections in this context a little.

\begin{definition}
  \label{def:weak-bigon}
  A \emph{weak bigon} between two curves~$C_1$ and~$C_2$ (intersecting
  transversally) is a
  map~$f:D^2\to\Sigma$ so that $\partial D^2$ is the union of two intervals~$A_1$
  and~$A_2$ so $f$ agrees with~$C_1$ on~$A_1$ and $f$ agrees
  with~$C_2$ on~$A_2$.
\end{definition}
%\nota{Still needs some cleanup.}

Hass and Scott~\cite{HS85:IntersectionCurves} have considered various
notions of ``bigons'' and whether they exist for various types of
embeddings with excess intersections; unfortunately, none of their
theorems are quite what we need.  Our notion of ``weak bigon'' is even
weaker than their weakest notion of bigon, since the boundary of the
bigon may wrap more than once around $C_1$ and/or $C_2$.

Note that the existence of a weak bigon is equivalent to a
lift~$\tilde C_1$ of~$C_1$ and a lift~$\tilde C_2$ of~$C_2$
intersecting more than once in the universal cover~$\tilde \Sigma$ of~$\Sigma$.

\begin{lemma}\label{lem:bigon-wrap}
  If $X$ is embedded and~$C$ is a loop, then no weak bigon
  between~$A$ and~$X$ wraps more than once around~$C$.
\end{lemma}

\begin{proof}
  Consider the annulus cover $\Sigma_C$ of $\Sigma$ corresponding to~$C$.  Then
  the lifts of~$C$ and~$X$ to~$\Sigma_C$ are embedded.  An arc bounding a
  bigon that wraps more than once around the lift of~$C$ is
  necessarily not embedded.
\end{proof}
%\nota{Give a simple picture here.}

\begin{lemma}\label{lem:taut-simple-bigon}
  Let $A \in \C_T(\Sigma)$ and $X \in \SC_S(\Sigma)$, with~$S$ and~$T$
  disjoint, and
  suppose we are given a transverse immersion of $X \cup A$.  Then this
  immersion has excess intersections between~$X$ and~$A$ if and only
  if there is a weak bigon between~$X$ and~$A$.
\end{lemma}

\begin{proof}
  If there is a bigon between~$X$ and~$A$, by
  Lemma~\ref{lem:bigon-wrap} it involves an arc of~$A$.  We can
  push that arc of~$A$ across the disk defining the bigon until it
  lies on the other side of~$A$.  When we do this
  we remove two intersections of~$A$ with~$X$ (at the two corners of
  the bigon), and, since~$X$ is simple, we do not introduce any new
  intersections between $A$ and $X$.

  Conversely, suppose that there is excess intersection between~$A$
  and~$X$.  By a result of Hass and
  Scott~\cite[Theorem~2.2]{HS94:ShorteningCurves}, there is a homotopy
  from $A \cup X$ to an immersion with minimal intersection which never
  increases the number of self-intersections.  If we put the homotopy
  in general position, this means that there is a sequence of $3\leftrightarrow3$,
  $2\to0$, $1\to0$, and $\partial 1\to0$ moves which makes $A\cup X$ taut.
%\nota{Need pictures of these moves.}
  Since there are excess intersections between~$A$ and~$X$, there will
  be at least one $2\to0$ move involving~$A$ and~$X$.  (This is the only
  possible move that reduces the number of $AX$ intersections.)
  Consider the configuration just before the first of such move; by
  definition, there will be an embedded bigon between $A$ and $X$.  We
  can follow this bigon backwards through the sequence of moves that
  preceded it to find a weak bigon in our original configuration.%\nota{This ``following backwards'' should be made clearer.  Perhaps
    %in universal cover?}
\end{proof}

\begin{question}
  Can excess intersections more generally be characterized by the
  existence of weak monogons, bigons, and boundary triangles?  The
  argument using curve shortening generalizes to show that if there is
  excess intersection, then there is also a weak monogon, bigon, or
  boundary triangle.  The other direction seems more difficult.
\end{question}
\DMGtodo[inline]{I am not sure if we want to leave this question here.}
\begin{proof}[Proof of Lemma~\ref{lem:smoothing}]
  By Corollary~\ref{cor:taut-extension}, we can suppose that~$X$
  and~$A$ are simultaneously embedded so that they are both taut.
  Then $i_X(A_+) \leq i_X(A)$, since the concrete configuration of $A_+$
  has $i_X(A)$ intersections with $X$, and a taut configuration can only
  have fewer intersections.  Similarly $i_X(A_-) \leq i_X(A)$.

  So we only need to show that equality is achieved at least once.
  Consider the lifts $\tilde A$ and $\tilde X$ of $A$ and $X$ to the
  universal cover of $\Sigma$, and a lift $\tilde x$ of $x$.  Name the 4
  rays of $\tilde A$ that meet at $\tilde x$ by $p$, $q$, $r$, and
  $s$, in cyclic order, so that $\tilde A_+$ joins $p$ to $q$ and $r$
  to $s$ and $\tilde A_-$ joins $p$ to $s$ and $q$ to $r$.
 % \nota{Picture here.}
  If $\tilde A_+ \cup \tilde X$ is not taut, then there is a bigon
  between the two curves.  (The other cases are ruled out, by
  assumption on $A$.)  Since the original $\tilde A$ was taut, this bigon
  must not have appeared originally; that is, it must have come from a
  triangle involving $X$ and two rays of $\tilde A$ from $\tilde x$; without
  loss of generality, these are $p$ and $q$.  Any strand of $\tilde X$ that
  crosses this triangle must meet both strands $p$ and $q$: it cannot
  meet one strand twice, by tautness of $\tilde A \cup \tilde X$, and it
  cannot meet the $\tilde X$ side, since $\tilde X$ is embedded.  Thus
  there is a unique innermost such triangle which involves the first
  intersections with $\tilde X$ along $p$ and $q$.  Similarly if
  $\tilde A_- \cup \tilde X$ is not taut, there is a triangle involving
  $X$ and (say) the first intersections along $q$ and $r$. But these
  two triangles together yield a $X$-$p$-$r$ bigon, contradicting the
  assumption that $A \cup X$ was taut.
\end{proof}